\documentclass[11pt,a4paper,xcolor=dvipsnames, leqno]{beamer}

\usetheme{show_and_tell}

\title[Introduction to Mathematical
Reasoning. Induction \& Contradiction Proofs]{An Introduction to Mathematical
Reasoning with Applications to
Induction and Contradiction Proofs}
\author[M\textsuperscript{a} Asunci\'on Jim\'enez Cordero]
{{{\bf M\textsuperscript{a} Asunci\'on Jim\'enez Cordero\\
{\href{mailto:asuncionjc@us.es}{\tt asuncionjc@us.es}}}}}
\institute{Show & Tell Karumi} % Opcional
\date{}


\defbeamertemplate*{title page}{customized}[1][]
{\hspace*{-0.3cm}
	\centering		
  	\begin{beamercolorbox}[wd=1.05\textwidth, ht = 1.4cm, rounded = true, shadow = true]{my_palette}
   		 \centering \usebeamerfont{title}\smallskip \inserttitle
   		 \smallskip
     \end{beamercolorbox}
   \insertdate\par
 
}

\begin{document}
%%%%%%%%%%%%%%%%%%%%%%%%%%%%%%%%%%%%%%%%%%%%%%%%%%%%%%%%%%%%%%%%%%%%%%%%%%%%%%%%
\begingroup 
    \setbeamertemplate{headline}{}
    \begin{frame}
        \titlepage
\begin{center}
\insertauthor\par
\end{center}
\begin{center}
        \scriptsize Departamento de Estad\'istica e Investigaci\'on Operativa\\ and Instituto de Matem\'aticas de la Universidad de Sevilla,\\
        Sevilla, Spain
        \end{center}
        
\vspace*{0.3cm}
\begin{center}
\includegraphics[scale=0.35]{imus_basico_oct.png}
\end{center}
\begin{center}
\scriptsize{Show \& Tell Karumi}
\end{center}
\begin{center}
\footnotesize{XX XXXX 2019}
\end{center}
\end{frame}
\endgroup

%%%%%%%%%%%%%%%%%%%%%%%%%%%%%%%%%%%%%%%%%%%%%%%%%%%%%%%%%%%%%%%%%%%%%%%%%%%%%%%%


\AtBeginSection[]{
\begin{without_headline}
  \begin{frame}{Outline}
    \tableofcontents[currentsection]
  \end{frame}
  \end{without_headline}
}


\begin{without_headline}
    \begin{frame}{Outline} % and our simple frame
        \tableofcontents
    \end{frame}
\end{without_headline}

%%%%%%%%%%%%%%%%%%%%%%%%%%%%%%%%%%%%%%%%%%%%%%%%%%%%%%%%%%%%%%%%%%%%%%%%%%%%%%%
\section{What is a proof?}
\begin{frame}
A \texthighlighted{theorem} is a statement that can be shown to be true.
\begin{itemize}
\item Hypothesis.
\item Thesis.
\end{itemize}
\begin{block}{Example}
\textbf{Bolzano's theorem:} If $f$ is a continuous function defined on a closed interval  $[a, b]$ such that $sign(f(a)) \neq sign(f(b))$.\\
Then there exists a point $c$ in the open interval $(a, b)$ satisfying $f(c) = 0$.
\begin{itemize}
\item Hypothesis:
\begin{itemize}
\item $f:[a, b]\rightarrow \mathbb{R}$ is continuous.
\item $sign(f(a)) \neq sign(f(b))$
\end{itemize}
\item Thesis:\\
\begin{itemize}
\item $\exists c\in (a, b): f(c) = 0$
\end{itemize}
\end{itemize}
\end{block}
\end{frame}
\section{Proof by contradiction}
\section{Proof by induction}
\section{Conclusions}


%%%%%%%%%%%%%%%%%%%%%%%%%%%%%%%%%%%%%%%%%%%%%%%%%%%%%%%%%%%%%%%%%%%%%%%%%%%%%%%
\begingroup 
    \setbeamertemplate{headline}{}
    \begin{frame}
        \titlepage
\begin{center}
\insertauthor\par
\end{center}
\begin{center}
        \scriptsize Departamento de Estad\'istica e Investigaci\'on Operativa\\ and Instituto de Matem\'aticas de la Universidad de Sevilla,\\
        Sevilla, Spain
        \end{center}
\begin{center}
\LARGE \texthighlighted{Thank you!}
\end{center}
\begin{center}
\includegraphics[scale=0.35]{imus_basico_oct.png}
\end{center}
\begin{center}
\scriptsize{Show \& Tell Karumi}
\end{center}
\begin{center}
\footnotesize{XX XXXX 2019}
\end{center}
\end{frame}
\endgroup

\end{document}
